\section{Each Element in Its Setting}
\label{sec:each-element}

The texts of this Mass come from two books with separate histories. The
antiphons and orations belong to the Missal's Twenty-eighth Week in Ordinary
Time and are prayed with the readings of all three years; the readings, psalm
and acclamation belong to the Lectionary's Year~A. Each element has its own
literary setting and its own textual facts, and the witnesses who expounded it
read it in that setting. What follows gives each element its context and the
witnesses' reading of it; the interpretations that follow this section carry
the arguments that join the elements.

\subsection{The question from the depths}

The Entrance Antiphon comes from Psalm 129 (Hebrew 130), one of the songs of
ascents, which the Vulgate heads \latin{Canticum graduum} and the
Douay--Rheims ``A gradual canticle''; the title names no author. The singer
cries ``out of the depths'' (v.~1), asks who could stand if the Lord marked
iniquities (v.~3), confesses that with him there is propitiation and that he
has waited for him by reason of his law (v.~4), and then turns to the people:
let Israel hope in the Lord, ``Because with the Lord there is mercy: and with
him plentiful redemption. And he shall redeem Israel from all his iniquities''
(vv.~7--8). The antiphon keeps the question and the first half of its answer
and closes with a vocative to the God of Israel, so that the name the psalm
speaks of at its end is spoken to at the antiphon's end.

Augustine hears the cry from the depths as the cry of every sinner, for ``this
mortal life is our deep''. Those who prosper in their sins are deepest of all:
``they are the more plunged in the deep, in proportion as they seem to be more
happy; for a deceitful happiness is itself a greater unhappiness.'' The
question of v.~3 is asked by one who saw ``that a clean heart trusting in its
own righteousness could not be found''. The answer turns on one word: ``And
what is this propitiation, except sacrifice? And what is sacrifice, save that
which hath been offered for us? The pouring forth of innocent blood blotted out
all the sins of the guilty'' (\work{Enarrationes in Psalmos} 129.1--3).
Robert Bellarmine reads the verb as a creditor's: to ``observe'' iniquities is
``to note down, to record, to make an entry, as a creditor would against a
debtor''. No one could pass that reckoning, because ``we, without his grace,
are not only unable to offer condign satisfaction, but we are even incapable of
seeing the enormity of the offence''; the sinner ``can transgress, but cannot
make satisfaction for the transgression, unless he be mercifully helped
thereto''. On v.~4 he takes the law by reason of which the psalmist waits to be
a law God laid on himself, ``to show no mercy to the impenitent, but to receive
the penitent'' (\work{Commentary on the Book of Psalms}, Ps~129:3--5).

\subsection{Grace that goes before and follows}

The Collect is short and exact. It asks that the Lord's grace \latin{semper et
praeveniat et sequatur}, always go before and follow, and that it keep the
assembly \latin{bonis operibus iugiter \ldots\ intentos}, constantly intent on
good works. The prayer is old. Its Latin stands in the Gregorian sacramentary
as H.~A. Wilson edited it in 1915, once among the prayers for vespers and
matins, where one manuscript has the Missal's order of words (p.~135), and
again among the Sunday prayers of the season after Pentecost (p.~174).

No Father comments on this prayer, but its two verbs have a precise
theological history. Thomas Aquinas asks whether grace is fittingly divided
into prevenient and subsequent, and answers from two psalm verses, ``his mercy
shall go before me'' and ``his mercy shall follow me'', the second of them the
sixth verse of this Sunday's psalm. Grace has five effects, he says: it heals
the soul, makes it will the good, makes it do the good, makes it persevere, and
brings it to glory; and each effect is prevenient with respect to the next and
subsequent with respect to the one before. He closes with a sentence of
Augustine's that he quotes: \latin{praevenit ut vocemur, subsequitur ut
glorificemur}, grace goes before that we may be called and follows that we may
be glorified (\work{Summa theologiae} I-II, q.~111, a.~3). He does not cite the
Collect. Two commentators on the Mass speak of the prayer itself. Lucien
Fromage, continuing Prosper Gu\'eranger's \work{Liturgical Year}, writes that
``unless grace prevent, that is, anticipate, us, we cannot have so much as the
thought of doing what is holy; and again, unless it follow up the inspirations
it has given us, and lead them to a happy termination, we shall never be able
to pass from the simple thought to the act of any virtue whatsoever'' (vol.~11,
p.~357). Blessed Ildefonso Schuster explains that ``Grace precedes the good
actions of our free will'', that ``the grace which has moved the will to the
act of volition accompanies, so to speak, this act and wholly pervades it'',
and that ``This truth of our holy Faith must make us very humble in the sight of
God'' (\work{The Sacramentary}, vol.~3, p.~142).

\subsection{A feast on the mountain}

The first reading belongs to the part of Isaiah that ends the oracles against
the nations with a vision of the whole earth judged and the Lord reigning.
Chapter 24 closes with ``the Lord of hosts'' reigning ``in mount Sion, and in
Jerusalem'' and glorified before his ancients (24:23). Chapter 25 opens with a
song of thanks to God, who has reduced ``the strong city to ruin'' and has been
``a strength to the poor, a strength to the needy in his distress'' (25:2, 4).
The appointed verses follow with ``And'': on ``this mountain'' the Lord makes
the feast for all people, destroys the bond and the web over the nations,
casts death down for ever, wipes away tears and takes away his people's
reproach; ``in that day'' they say, ``Lo, this is our God, we have waited for
him''; for the hand of the Lord rests on this mountain (25:6--10a). The rest of
v.~10 and vv.~11--12, on Moab trodden down, lie outside the reading.

The Latin forms differ in ways that matter to the reception. The \work{Ordo}'s
title, \latin{absterget lacrimam}, follows the \work{Nova Vulgata}, where the
Clementine Vulgate reads \latin{auferet}; the \work{Nova Vulgata} speaks in v.~6
of pure and refined wine where the Clementine speaks of the vintage. Both read
in v.~8 \latin{Praecipitabit mortem in sempiternum}, and the Douay--Rheims
follows them: ``He shall cast death down headlong for ever.'' Paul quotes the
same verse in another form, ``Death is swallowed up in victory'' (1~Cor 15:54).
Jerome, expounding the chapter, sets the Septuagint beside the Vulgate, and
says of the Greek translators' addition, \latin{Trade omnia haec gentibus},
``give all these to the nations'', that they gave \latin{non Scripturae verba,
sed suum sensum}, not the words of Scripture but their own sense of them
(\work{Commentarii in Isaiam} VIII, PL 24, col.~300). No Father cited here read
the Lectionary's Latin, and none of their lemmas is the Lectionary's text.

Jerome reads the banquet as the feast the Lord makes after the Passion for all
the nations on Mount Sion, ending when death is swallowed up and Christ's
kingdom comes; his sentence, and the contrast it draws, is given in full under
the first interpretation. On v.~9 he gives the waiting a precise content:
\latin{Ecce Deus noster iste, quem increduli hominem tantum putabant; et
exspectavimus eum, hoc est, verbis eius credidimus, quia sua promissa complebit,
et salvabit nos}: this is our God, whom unbelievers thought only a man; and we
waited for him, that is, we believed his words, because he will fulfil his
promises and save us (cols.~300--301). Thomas Aquinas reads the chapter as
thanksgiving, and vv.~6--10 as the exaltation of the good. He is the one
witness cited here who joins this reading to the Gospel: the feast is made
\latin{in monte hoc, quia ibi Christus passus est, unde nobis omnia bona
provenerunt; vel quia ibi iudicabit, ut dicitur Matth. 22: Ecce prandium meum
paravi}, on this mountain because Christ suffered there, from which all good
things have come to us, or because he will judge there, as Matthew 22 says,
``Behold, I have prepared my dinner'' (\work{Expositio super Isaiam} c.~25,
Parma, pp.~501--502). He reads the casting down of death with 1~Corinthians
15:54 and Hosea 13:14 and the wiping of tears with Apocalypse 21:4, and he
notes that some referred the passage to the judgment shown in the death of
Holofernes. The third \work{Mystagogical Catechesis}, transmitted under the name
of Cyril of Jerusalem, quotes the verse in its Septuagint form to the newly
anointed: ``on this mountain shall the Lord make unto all nations a feast; they
shall drink wine, they shall drink gladness, they shall anoint themselves with
ointment''; the mountain, it says, is the Church (3.7).

\subsection{The shepherd's table}

The psalm is the twenty-second in the Vulgate's count, headed \latin{Psalmus
David}, ``A psalm for David''. The Latin says that the Lord ``ruleth'' the
singer, where the Hebrew says that he shepherds him; the images that follow
are a shepherd's, pasture and water, the right paths, the rod and staff in the
shadow of death, and then a host's, a table prepared in the sight of enemies,
the head anointed, the cup filled, and mercy following the guest until he
dwells in the house of the Lord. The response is taken from the second half of
the last verse.

Augustine hears the whole psalm as the Church speaking to Christ: ``The Lord
Jesus Christ is my Shepherd, `and I shall lack nothing.'\,'' The water of
refreshment is baptism; the paths of justice are ``the narrow ways, wherein few
walk''; the shadow of death is ``this life''. The table is prepared ``that I
should no more be fed as a babe with milk, but being older should take meat'';
the oil is spiritual joy and the cup is ``yielding forgetfulness of former vain
delights''. The last verse reaches past this life: ``Thy mercy shall follow me
not here only, but also that I may dwell in the house of the Lord for ever''
(\work{Enarrationes in Psalmos} 22.1--6). Aquinas, in his commentary on the
Psalter, says that the psalm can signify the people's return from Babylon and
signifies Christ's return from the world to heaven, and he hears it spoken
\latin{in persona Ecclesiae} of Christ the shepherd. ``I shall want nothing''
is said of what is necessary for salvation, and of the future: \latin{in futuro
omnem sufficientiam habebimus, quia nihil deerit nobis, quoniam habebimus
Deum}, in the future we shall have every sufficiency, because nothing will be
lacking to us, since we shall have God. The table is the teaching of
Scripture, \latin{vel mensam sacramentalem, scilicet altaris}: the table of the
altar, for the table is threefold, of the old law, of the New Testament, and
\latin{in patria}, in the homeland. The cup is the gift of divine love that
inebriates, or the blood of Christ; and the house of the Lord is the Church
\latin{per gratiam} and heaven \latin{per gloriam} (\work{Postilla super
Psalmos} 22). Robert Bellarmine finds the sheep's pasture in ``the knowledge of
God, his sacraments, especially the Eucharist''; the table in the presence of
enemies means that ``God provides great consolations to meet great
tribulations'', pictured ``as if we were enjoying a feast''; and the mercy that
follows is ``the last good, that brings to the supreme good'' (\work{Commentary
on the Book of Psalms}, Ps~22:1--6). The fourth \work{Mystagogical Catechesis}
calls the table ``that mystical and spiritual Table, which God hath prepared for
us'' and the cup the one which Jesus took in his hands, saying ``This is My
blood'' (4.7).

\subsection{Content in want and in plenty}

The second reading is the last of four Sundays on which Year~A reads
Philippians, and it comes from the letter's close. Paul has thanked the
Philippians for their renewed care for him (4:10) and said that he has
``learned, in whatsoever state I am, to be content therewith'' (4:11). The
appointed verses give the content of that learning and its source: to be
brought low and to abound, to be full and to be hungry, and the strength of
the one who strengthens him; and they commend the Philippians for sharing his
tribulation. The verses left out recall that they alone gave to him from the
beginning of the Gospel, that he seeks ``not the gift'' but ``the fruit'', and
that he is filled by what Epaphroditus brought, ``an odour of sweetness, an
acceptable sacrifice, pleasing to God'' (4:15--18). The reading resumes with
God's supplying of all their want and closes with the doxology. The
\work{Ordo}'s title for the reading is the verse \latin{Omnia possum in eo, qui
me confortat}.

John Chrysostom and Thomas Aquinas both read Paul's contentment as a virtue
that has to be learned and is needed as much in plenty as in want, and both
read its strength as Christ's; their words are given under the third
interpretation. On v.~19 the two read differently because they read different
texts. Chrysostom's text is a blessing: ``Behold how he invokes blessings upon
them, as poor men do''; and he reports that copies read ``every need'', ``every
grace'' or ``every joy'' (\work{Homilies on Philippians} 15). Aquinas expounds
the Vulgate's \latin{impleat omne desiderium}: first as God's return to the
givers who had filled Paul's need, and then of desire filled in glory (\work{Super Philippenses} c.~4, lect.~2,
Dessain, p.~403). The third \work{Mystagogical Catechesis} gives v.~13 to the
newly anointed, who, having put on ``the whole armour of the Holy Ghost'', are
to stand against the adversary ``saying, I can do all things through Christ
which strengtheneth me'' (3.4).

\subsection{The hope of our calling}

The Alleluia verse turns Paul's prayer for the Ephesians into the assembly's
prayer for itself. In the letter Paul, having heard of their faith and love,
prays that the Father of glory give them the spirit of wisdom and revelation,
the eyes of their heart enlightened to know the hope of his calling, the riches
of the glory of his inheritance in the saints and the greatness of his power
toward those who believe (Eph 1:15--19). The verse keeps the enlightening and
the hope and makes the calling ``ours''. John Chrysostom says that Paul wants
them to understand ``to what blessings they are called'', and of the hope he
says, ``It is as yet, he means, hidden, but not so to the faithful.'' The power
that brought them to believe is greater than the power that raises the dead:
``to persuade souls, is a thing far more miraculous than to raise a dead
body'', because ``He would thus have us become good of our own will''; and
``there is need of grace in order that the understanding may perceive it''
(\work{Homilies on Ephesians} 3). Aquinas reads the enlightened eyes as the gift
of wisdom, given for the knowledge of God, \latin{contra eos, qui habent oculos
illuminatos tantum ad temporalia cognoscenda}, against those whose eyes are
enlightened only to know temporal things, with Jeremiah's \latin{non glorietur
dives in divitiis suis}, let not the rich man glory in his riches (Jer 9:23).
The hope of the calling is \latin{virtus spei, et de quanta re sit}, the virtue
of hope and the greatness of the thing hoped for (\work{Super Ephesios} c.~1,
lect.~6, Dessain, p.~272).

\subsection{The king's wedding}

The parable is the third that Jesus tells in the temple after the chief
priests and the ancients of the people have asked by what authority he acts
(Mt 21:23). The parable of the vineyard has just ended with the saying that the
kingdom will be taken from its tenants and given to a nation yielding its
fruits, and ``the chief priests and Pharisees \ldots\ knew that he spoke of
them'' and feared the crowds (21:43--46). Matthew continues, ``And Jesus
answering, spoke again in parables to them'' (22:1). The parable that follows
ends the series; after it the Pharisees take counsel how to ensnare him in his
speech, and the question of tribute begins (22:15), the Gospel of the following
Sunday. The \work{Ordo}'s incipit names the chief priests and elders of the
people as those addressed, the words of 21:23.

The story moves in three steps. The king's invitation is refused twice, the
second time in spite of the word that ``all things are ready''; some neglect
it and go off to farm and merchandise, the rest kill the servants, and the
king's armies destroy the murderers and burn their city (22:2--7). The
invitation then goes out to the roads: the Latin has \latin{ad exitus viarum},
where the roads go out, which the Douay--Rheims renders ``into the highways'',
and the servants gather ``both bad and good'' (22:8--10). In the longer form
the king comes in, finds a guest without a wedding garment, and has him bound
and cast out, and the saying on the called and the chosen closes the parable
(22:11--14). The shorter form ends with the hall full. Under no.~80 of the
\work{Praenotanda} the choice between the forms is pastoral, made for the
hearers' capacity to listen with profit; nothing in the books makes either the
proper one.

Jerome notices how the parable names the host: when he invites and does works
of clemency he is \latin{homini regi}, a man who is king, but when he comes to
vengeance \latin{homo siletur, et rex tantum dicitur}, the man is left unsaid
and he is called king alone (\work{Commentarii in Matthaeum} III, PL 26,
col.~160). Aquinas, giving Origen's reason, explains the same phrase
differently: he is called a man who
is king because he now rules us in a human way, \latin{Sed cum videbitur sicuti
est, tunc erit Rex, quia tunc secundum se reget}, but when he is seen as he is,
then he will be King, because then he will rule according to himself
(\work{Super Matthaeum} c.~22, Venice, p.~279). Gregory the Great and Aquinas
both ask whether Luke's great supper (Lk 14:16--24) is the same parable, and
both hold that it is another: for Gregory, Matthew's wedding signifies the
present Church, where some enter who will go out again, and Luke's supper the
eternal and last banquet, which no one who has entered will leave
(\work{Homiliae in Evangelia} 38.1). Gregory yields the point, \latin{salva
fide}, to anyone who holds the two to be one.

The text itself is stable. The editions reported in the apparatus of the SBL
Greek New Testament differ only in small ways: at v.~7 the Byzantine text, with
the Clementine Vulgate and the Douay--Rheims, has the king hear before he is
angry, which the critical texts and the \work{Nova Vulgata} do not; at v.~10
one edition reads the bridal hall for the wedding; at v.~13 the Byzantine text
adds ``take him away''. None of these changes who is invited, who refuses, who
is gathered or why the guest is cast out. Chrysostom's English text says that
the invited ``made light of it'' where the Vulgate says \latin{neglexerunt}.

What the parable's details mean is the matter of the three interpretations
that follow. Two disputes among the Fathers run through them and are stated
here once. The first is where the wedding stands: Gregory places it in the
present Church, Hilary of Poitiers makes it the sacrament of eternal glory to
be received at the resurrection, Augustine speaks of two feasts, and Aquinas
gives both, following Gregory for the parable. The second is the burnt city of
v.~7: John Chrysostom reads it of the destruction of Jerusalem under Vespasian
and Titus; Jerome and Aquinas give that reading and another; Gregory and
Hilary read the armies as angels and the fire as the eternal punishment of the
persecutors. Both disputes are developed under the first interpretation, and a
third, on what the wedding garment is, under the second.

\subsection{Gifts offered for glory}

The Prayer over the Offerings asks the Lord to receive the faithful's prayers
together with the sacrificial offerings, so that through these acts of devout
worship they may pass over to heavenly glory. Its body stands word for word in
Wilson's edition of the Gregorian sacramentary as the prayer over the
offerings for the station at Saint Paul's on the Tuesday of Easter week
(p.~61), and again on a Sunday after the octave of Easter (p.~167). No
commentary on the prayer was found; what the
interpretations say of it is their own reading of its words, \latin{per haec
piae devotionis officia} and \latin{ad caelestem gloriam}.

\subsection{The rich who hungered}

The first Communion antiphon comes from Psalm 33 (Hebrew 34), whose title in
the Vulgate and the Douay--Rheims ties it to David's changing his countenance
before Achimelech, the episode the Douay--Rheims margin locates in 1~Kings 21.
It is the psalm of ``O taste, and see that the Lord is sweet'' (v.~9). Verse 10
promises that there is no want to those who fear the Lord, and v.~11, the
antiphon's verse, sets against them the rich who ``have wanted, and have
suffered hunger''; v.~12 invites children to learn the fear of the Lord. The
rich of the verse are the Latin and Greek reading; the Hebrew has young lions.

Augustine meets the obvious objection, that the wicked rich often die in
plenty, and answers it by the spiritual sense, closing his comment by joining
the verse himself to the Gospel of John: the bread the rich lack is the Living
Bread (\work{Enarrationes in Psalmos} 33, sections 13--14 of the English
text). Robert Bellarmine gives the same ``higher meaning''. Both are quoted
under the third interpretation, where their reading carries most weight.

\subsection{Seeing him as he is}

The second Communion antiphon comes from the First Letter of John: ``Behold
what manner of charity the Father hath bestowed upon us, that we should be
called and should be the sons of God'' (3:1). The sons do not yet know what
they shall be; they know that when he appears they shall be like him, because
they shall see him as he is (3:2); and everyone who has this hope in him
sanctifies himself, as he is holy (3:3). The antiphon gives only the knowledge
and supplies its subject, \latin{Dominus}, which the letter leaves unexpressed.

Augustine asks what it is to see him ``as He is''. The form of the servant
even the wicked will see at the judgment, but to see Christ in the form of God
``is to the bad impossible''; it is a vision ``surpassing all things: because
from it are all things beautiful'' (\work{Homilies on the First Epistle of
John} 4.5). Aquinas makes the verse the authority for his article that the
blessed see the essence of God (\work{Summa theologiae} I, q.~12, a.~1), and he
explains that ``as'' determines what is seen, not how fully it is seen: ``we
shall see Him to be as He is'', but one will see more perfectly than another
(a.~6, ad~1). What makes the difference is charity: ``he will have a fuller
participation of the light of glory who has more charity; because where there
is the greater charity, there is the more desire; and desire in a certain
degree makes the one desiring apt and prepared to receive the object desired''
(a.~6).

\subsection{Fed, and made partakers}

The Prayer after Communion asks the divine majesty that, as he feeds the
assembly with the Body and Blood of Christ, he make them partakers of the
divine nature. An older form of it stands in the July section of the Verona
sacramentary, without the vocative \latin{Domine} that the Missal adds; that
witness was read only in a machine-read text of Feltoe's edition (1896). Its
closing words are those of 2~Peter 1:4, where the promises given
through the knowledge of him who called us are given ``that by these you may
be made partakers of the divine nature: flying the corruption of that
concupiscence which is in the world''. No commentary on the prayer was found,
but the joining of its two halves has a witness. The fourth \work{Mystagogical
Catechesis} tells the newly baptized that ``by partaking of the Body and Blood
of Christ'' they are made ``of the same body and the same blood with Him'', and
``thus it is that, according to the blessed Peter, we become partakers of the
divine nature'' (4.3). Aquinas, without reference to the prayer, calls grace
itself \latin{participatio divinae naturae} and derives the infused virtues
from it (\work{Summa theologiae} I-II, q.~110, a.~3), and he says that God
alone deifies, \latin{communicando consortium divinae naturae}, by
communicating a share in the divine nature (q.~112, a.~1).
